% Isolated after the CP34 canonical cutoff; no canonical or shared-input mutation.
\subsection{A determinant wedge and the degree of the regular inertia wall}

Work over an algebraically closed field of characteristic two. Let
\(\eta:A\to B\) be an etale isogeny of positive odd degree, with
\((B,\Theta_B)\) nonordinary and principally polarized, and put
\[
 L=\eta^*\mathcal O_B(3\Theta_B),\qquad
 N=h^0(A,L),\qquad X=\mathbf P(H^0(A,L)),\qquad r=N-1.
 \tag{DE-WP1}
\]
Write \((d_1,d_2)\), with \(d_1\mid d_2\), for the polarization
type of \(L\). Thus \(N=d_1d_2\), both integers are odd, and
\(d_1\ge3\): the line \(L\) is the cube of the ample line
\(\eta^*\mathcal O_B(\Theta_B)\).
The whole-coefficient-space Euler model (OW12) is the rank-\(r\)
vector bundle \(\mathcal E=\Omega_X^1\). Its actual coherent
Frobenius is a regular map
\[
 \Phi:F_X^*\mathcal E\longrightarrow\mathcal E,
 \tag{DE-WP2}
\]
where \(F_X\) is absolute Frobenius. This is the linearization of
the semilinear coherent map; relative Cartier on the abelian surface
is still the operator used to construct its coefficient matrices.
Generic sharpness makes \(\det\Phi\ne0\). Its divisor \(D\)
has degree \(N\), because
\(\det\mathcal E=\mathcal O_X(-N)\) and hence
\(\det\Phi\) is a section of \(\mathcal O_X(N)\).

Let \(G\) be the geometric relative Lang image and let
\(K\subset V\) be its canonical proper submodule from the full-radical
structure theorem. Put \(k=\dim_{\mathbf F_2}K\) and
\(m=\dim(V/K)=r-k\). No extension of \(K\) across every component
of \(D\) is assumed below.

\begin{lemma}[Divisorial poles of an invariant determinant wedge]
\label{lem:DE-Lang-wedge-pole-bound}
\status{Candidate; complete proof supplied for independent audit.}
If \(k>0\), the determinant of \(K\) gives a nonzero rational
section
\(w\) of \(\bigwedge^k\mathcal E\) satisfying
\[
                 (\bigwedge^k\Phi)(F_X^*w)=w.
 \tag{DE-WP3}
\]
At any prime divisor \(Z\subset X\), its pole order is at most
\(\operatorname{ord}_Z(\det\Phi)\).
If \(Z\) is the closure of a component of the common maximal-rank
wall \(H_{\max}\) in (OW9), then \(w\) is regular along \(Z\).
\end{lemma}
\begin{proof}
On the generic Frobenius-bijective locus, the Lang local system is
the fixed-vector sheaf of \(\Phi\) in \(\mathcal E\).
A basis of \(K\) over \(\mathbf F_2\) consists of linearly
independent vectors over the algebraically closed function field
as well: a Frobenius-bijective semilinear space is the scalar
extension of its Lang fixed space. Their wedge is nonzero.
Every change of this basis has determinant in
\(\mathbf F_2^\times=\{1\}\), so the wedge is independent of
the basis and invariant under the geometric fundamental group.
It therefore descends to a rational section on \(X\), and its
fixed-vector equation is exactly (DE-WP3).

At the generic point of \(Z\), trivialize \(\mathcal E\) over
its discrete valuation ring. Put \(a=\operatorname{ord}_Z\det\Phi\).
The matrix of \(\Phi\) is integral and generically invertible.
Its Smith exponents \(a_1,\ldots,a_r\) are nonnegative and have
sum \(a\). Every entry of
\((\bigwedge^k\Phi)^{-1}\) has valuation at least \(-a\): its
largest possible denominator exponent is a sum of \(k\) of the
\(a_i\), which is at most their full sum. If \(v(w)\) denotes
the minimum valuation of its wedge coordinates, (DE-WP3) gives
\[
       2v(w)=v(F_X^*w)
             \ge v(w)-a,
       \qquad v(w)\ge-a.
 \tag{DE-WP4}
\]
The equality on the left uses the actual absolute Frobenius
pullback of those coordinates. This proves the pole bound without
assuming the individual entries of \(w\) are regular.

At a geometric generic point of a component of \(H_{\max}\),
Theorem~\ref{thm:odd-cover-full-hyperplane-inertia} gives stable
rank \(r-1\) and a full hyperplane radical in inertia. The
regular Lang vectors form its fixed hyperplane: their unique
etale lifts exist over a strictly henselian local ring because
the Lang equations have fibre Jacobian \(-I\), and there are
exactly \(2^{r-1}\) of them. This is the same regular-part
comparison used in Proposition~\ref{prop:DE-Lang-deck-slice-quotient-eight}.

The canonical \(K\) lies in the fixed hyperplane of every full
radical contained in \(G\). Indeed applying the full-radical
structure theorem to any such radical gives the same unique
maximal proper submodule. Thus all its vectors, and their
determinant wedge, are regular in this local comparison.
Regularity descends to the original divisorial valuation.
Only this rank-\((r-1)\) regular part is used; the full generic
Lang space is not specialized through a rank drop.
\end{proof}

\begin{lemma}[A theta-invariant divisor has a constrained degree]
\label{lem:DE-theta-invariant-divisor-degree}
\status{Candidate; complete proof supplied for independent audit.}
Let \(L\) be an ample line bundle of odd polarization type
\((d_1,d_2)\) on an abelian surface. Every nonzero effective
divisor on \(\mathbf P(H^0(L))\) invariant under the projective
theta-group action has degree divisible by \(d_2\).
\end{lemma}
\begin{proof}
Let \(f\) be its homogeneous equation, of degree \(e\).
Its line \(kf\) is stable under the theta group \(\mathcal G(L)\).
Consequently every pair of theta operators commute on this
one-dimensional line. The theta commutator pairing on \(K(L)\)
is perfect; this is the full nondegeneracy statement in
\citep[Proposition~8.19 and Corollary~8.20, p.~119]
{EdixhovenVanDerGeerMoonenAV}.
Since the group \(K(L)\) has exponent \(d_2\), choose an
element of order \(d_2\). Its paired character also has exact
order \(d_2\), so some commutator of theta operators is a
primitive \(d_2\)-th root of unity \(\zeta\) acting scalarly
on \(H^0(L)\). On degree-\(e\) equations it acts by
\(\zeta^{-e}\), which must be one on \(kf\). Hence
\(d_2\mid e\). The opposite convention for the action on
equations changes the sign of the exponent and has the same
conclusion. Oddness ensures that these roots of unity are
reduced and distinct; no inseparable character argument is used.
\end{proof}

\begin{theorem}[The radical-generated quotient has dimension at least \(d_2\)]
\label{thm:DE-Lang-theta-wedge-quotient-bound}
\status{Candidate; complete proof supplied for independent audit.}
In (DE-WP1), let \(D_{\mathrm{good}}\) be the reduced union
of the closures in \(X\) of all components of \(H_{\max}\).
It is a nonzero effective divisor of degree \(e\), and
\[
             d_2\mid e,\qquad 0<e\le N.
 \tag{DE-WP5}
\]
If \(K\ne0\), then
\[
             k\le N-e-1,
             \qquad m\ge e\ge d_2.
 \tag{DE-WP6}
\]
If \(e=N\), then \(K=0\) and the image is full.
In all cases, including \(K=0\), one has \(m\ge d_2\).
\end{theorem}
\begin{proof}
The common wall is nonempty by
Proposition~\ref{prop:odd-cover-common-rank-wall-nonempty}.
Every component is a smooth reduced divisor in the common
maximal-rank open, and the actual Euler determinant has order
one there by Theorem~\ref{thm:odd-cover-full-hyperplane-inertia}.
Its distinct closures therefore define a reduced nonzero divisor
\(D_{\mathrm{good}}\le D\), proving \(0<e\le N\).

This divisor is invariant under the projective theta group of
\(L\). To see this on the actual maps, a theta operator is a
translation of \(A\) followed by a line-bundle isomorphism.
It intertwines relative Cartier, with the corresponding
Frobenius-twisted operator on its target; invariant volume forms
are unchanged by translation. Hence it preserves both ranks
defining \(U_{\max}\), and also preserves the actual Euler
determinant wall. It consequently permutes the components of
\(H_{\max}\) and their closures. This is invariance of the
whole reduced union; no individual component is presumed fixed.
The preceding lemma gives \(d_2\mid e\).

Suppose \(k>0\). Set \(D_{\mathrm{bad}}=D-D_{\mathrm{good}}\),
an effective divisor of degree \(N-e\). The wedge lemma makes
\(w\) regular along every component removed from \(D\), and
bounds its poles by the remaining determinant multiplicities
at every other divisor. Since \(X\) is smooth and the wedge
bundle is locally free, these divisorial bounds give a nonzero
global section
\[
  w\in H^0\!\left(X,\bigwedge^k\Omega_X^1
                                  (D_{\mathrm{bad}})\right)
    =H^0\!\left(\mathbf P^{N-1},\Omega^k(N-e)\right).
 \tag{DE-WP7}
\]
The extension across codimension at least two uses only
reflexivity of a locally free sheaf on the regular variety.

For completeness, \(H^0(\mathbf P^{N-1},\Omega^k(t))=0\)
for \(1\le k\le N-1\) and \(t\le k\) in characteristic
two as well. The exterior Euler sequence embeds
\(\Omega^k(t)\) into
\(\bigwedge^k H^0(\mathcal O(1))\otimes\mathcal O(t-k)\).
For \(t<k\) the latter has no sections. For \(t=k\), a
section would be a constant exterior tensor whose contraction
with the universal coordinate vector is identically zero.
Independence of the coordinates makes contraction by every
basis vector zero, so that tensor is zero. This proves the
vanishing without using a characteristic-zero Bott theorem.

Equation (DE-WP7) therefore forces \(N-e\ge k+1\), proving
(DE-WP6). If \(e=N\), it rules out every \(k>0\); the
full-radical theorem then gives the full group. If \(k=0\)
from the outset, \(m=N-1\ge d_2\), since \(d_1\ge3\).
This establishes the unconditional last assertion without
applying the positive-wedge argument to a zeroth exterior power.
\end{proof}

\begin{corollary}[Cyclic and principal cubic systems]
\label{cor:DE-Lang-wedge-cyclic-principal-bounds}
\status{Candidate; no new full-image degree asserted.}
For a cyclic odd pullback of degree \(\ell\) of a cubic principal
system, with \(\ell\) an odd prime, \(L\) has type
\((3,3\ell)\), and
\[
          \dim(V/K)\ge3\ell,\qquad \dim K\le6\ell-1.
 \tag{DE-WP8}
\]
The degree \(e\) is one of \(3\ell,6\ell,9\ell\);
in the last case the image is full.
For every fixed nonordinary principally polarized surface and
every positive odd \(d\), the system \(|3d\Theta|\) satisfies
\[
    \dim(V/K)\ge\max\{8,3d\},\qquad
    \dim K\le9d^2-1-\max\{8,3d\}.
 \tag{DE-WP9}
\]
Its main nonzero Lang component has degree at least
\[
       2^{9d^2-1}-2^{\,9d^2-1-\max\{8,3d\}}.
 \tag{DE-WP10}
\]
\end{corollary}
\begin{proof}
For the cyclic source the underlying pulled-back principal line
has type \((1,\ell)\), so its cube has the displayed type.
Apply the theorem, with \(N=9\ell\).
For the fixed principal system use its exact outgoing odd etale
cubic-principal presentation. Its type is \((3d,3d)\).
Combine the theorem with the unconditional eight-dimensional
deck-slice bound. In this fixed-principal scope the main-orbit
lemma gives the exact main degree \(2^{9d^2-1}-2^{\dim K}\),
which implies (DE-WP10). At \(d=1\) the bound gives \(K=0\),
consistent with the already proved degree-one result.
\end{proof}

\begin{remark}[Poles at the remaining wall are retained]
\label{rem:DE-Lang-wedge-pole-scope}
\status{Candidate scope; a stronger bound, not a general vanishing.}
The proof does not extend \(K\), its vectors or their wedge
regularly across every determinant divisor. It permits poles
up to the full determinant multiplicity on the remaining wall.
Even multiplicity alone is not used to infer trivial inertia.
The theta-degree constraint and exterior-Euler vanishing give
the stated numerical bound despite those possible poles.
In particular the cyclic-source full-image hypotheses in the
compatible-intermediate criterion remain open. No new exact
higher degree or unrestricted A.4 conclusion is asserted.

All field and bundle maps are the actual Euler-cohomology
maps on the entire coefficient space. The only additional
source fact is theta-group nondegeneracy; its full original
proof in Proposition~8.19 and Corollary~8.20 was reread in the
retained text. The pole and exterior-Euler arguments are
written out above. No numerical experiment is needed or claimed.
\end{remark}
